% !TeX program = lualatex
% ================================================================
% Urdu Parallel Text Version of PercentagesQuickQuestions
%
% Generated by the server using the following settings:
%
%   language            Urdu (urd)
%   text direction      right to left
%   font                Noto Nastaliq Urdu
%   fallback font       Noto Serif
%   built from          latex/quickQuestions/percentagesQuickQuestions.tex
%   vocabulary key      latex/quickQuestions/percentagesQuickQuestions/L2/urd/percentagesQuickQuestions_urduKey.tex
%   tier 3 vocabulary   green   RGB 0 128 0
%   English text        black   RGB 0 0 0
%   Urdu text           purple  RGB 112 48 160
%   layout macros       version 3
%
% Please don't edit any information in this comment block - the server
% compares it against the current settings and rebuilds this file when
% they differ, so changes here would be overwritten. However, feel free
% to make updates to the LaTeX below if you want to edit it.
% ================================================================

\documentclass{beamer}
% ================================================================
% Copied from the English sheet this was translated from, so that
% anything it sets up for itself still works here
% ================================================================
\usepackage{amsmath}
\usepackage{amssymb}
\usepackage{tikz}
\usetikzlibrary{decorations.pathreplacing}

% ================================================================
% Quick Questions deck: percentages.
%
% Every question gets two slides: the question on its own for the
% mini whiteboards, then the same question again with the solution
% underneath in red. Within each topic the ten questions get
% gradually harder.
% ================================================================

\setbeamertemplate{navigation symbols}{}
\setbeamertemplate{footline}{}

% Topic divider.
\newcommand{\qtopic}[1]{%
  \section{#1}%
  \begin{frame}[plain]
    \vfill
    \begin{center}
      {\usebeamercolor[fg]{structure}\Huge\bfseries #1}
    \end{center}
    \vfill
  \end{frame}%
}

% \qq[<question size>][<solution size>]{<question>}{<solution>}
% The two optional sizes drop the type for the wordier questions and
% the longer solutions, so nothing runs off the slide.
\NewDocumentCommand{\qq}{O{\Huge}O{\LARGE}+m+m}{%
  \begin{frame}
    \vfill
    \begin{center}
      \begin{minipage}{0.92\textwidth}
        \centering #1 #3
      \end{minipage}
    \end{center}
    \vfill
  \end{frame}%
  \begin{frame}
    \vfill
    \begin{center}
      \begin{minipage}{0.92\textwidth}
        \centering #1 #3
        \par\vspace{0.9em}
        {\color{red}#2 #4}
      \end{minipage}
    \end{center}
    \vfill
  \end{frame}%
}

% Stepped working, written in the form
%   80% is 64 / 10% is 8 / 100% is 80
% Rows are separated with \ and the two columns with &.
\newcommand{\steps}[1]{%
  \begin{tabular}{@{}r@{\ is\ }l@{}}#1\end{tabular}%
}

% \barmodel{<cells>}{<cells spanned by the lower brace>}{<lower label>}{<upper label>}
% The whole bar is 9cm wide and represents the upper label; the lower
% brace picks out the part of it named by the lower label.
\newcommand{\barmodel}[4]{%
  \begin{tikzpicture}[x=1cm,y=1cm]
    \pgfmathsetmacro{\cw}{9/#1}
    \fill[red!12] (0,0) rectangle ({#2*\cw},0.7);
    \foreach \i in {1,...,#1}
      {\draw[red!60!black] ({(\i-1)*\cw},0) rectangle ({\i*\cw},0.7);}
    \draw[red!60!black,thick] (0,0) rectangle (9,0.7);
    \draw[red,decorate,decoration={brace,amplitude=4pt}] (0,0.9) -- (9,0.9)
      node[midway,above=4pt,red,font=\normalsize]{#4};
    \draw[red,decorate,decoration={brace,mirror,amplitude=4pt}] (0,-0.2) -- ({#2*\cw},-0.2)
      node[midway,below=4pt,red,font=\normalsize]{#3};
  \end{tikzpicture}%
}

% Contents-slide bullets, colour-coded by difficulty band:
% green for the first two topics, orange for the next two,
% red for the last three.
\setbeamertemplate{section in toc}{%
  \ifnum\inserttocsectionnumber<3
    \textcolor{green!55!black}{$\bullet$}%
  \else\ifnum\inserttocsectionnumber<5
    \textcolor{orange}{$\bullet$}%
  \else
    \textcolor{red}{$\bullet$}%
  \fi\fi
  ~\inserttocsection%
}

\title{Percentages: Quick Questions}
\subtitle{Mini whiteboard questions, with solutions}
\author{}
\date{}

% Characters this language's font has not got are borrowed from another,
% rather than being silently dropped - see docs/translated-sheet-latex.md
\directlua{%
  if luaotfload and luaotfload.add_fallback then
    luaotfload.add_fallback("eallatinfallback",
      { "Noto Serif:mode=harf;" })
  else
    tex.error("this luaotfload is too old to register a fallback font")
  end
}
\usepackage[english,bidi=basic]{babel}
\babelprovide[import]{urdu}
\babelfont[urdu]{rm}[Renderer=Harfbuzz, BoldFont={Noto Nastaliq Urdu}, ItalicFont={Noto Nastaliq Urdu}, BoldItalicFont={Noto Nastaliq Urdu}, RawFeature={fallback=eallatinfallback}]{Noto Nastaliq Urdu}
\babelfont[urdu]{sf}[Renderer=Harfbuzz, BoldFont={Noto Nastaliq Urdu}, ItalicFont={Noto Nastaliq Urdu}, BoldItalicFont={Noto Nastaliq Urdu}, RawFeature={fallback=eallatinfallback}]{Noto Nastaliq Urdu}
\newcommand{\ealtext}[1]{\foreignlanguage{urdu}{#1}}
\newcommand{\ealtextblock}[1]{{\raggedleft\foreignlanguage{urdu}{#1}\par}}
% ================================================================
% EAL layout helpers
% ================================================================
\usepackage{xcolor}

\definecolor{ealkeycolour}{RGB}{0,128,0}
\definecolor{ealenglish}{RGB}{0,0,0}
\definecolor{ealtwocolour}{RGB}{112,48,160}

\newsavebox{\ealboxen}
\newsavebox{\ealboxtr}
\newlength{\ealglosswd}
\newlength{\ealglossraise}
\setlength{\ealglossraise}{1.15em}

% a tier 3 word, in English and in translation alike
\newcommand{\ealkey}[1]{\textcolor{ealkeycolour}{#1}}
\newcommand{\ealkeytr}[1]{\textcolor{ealkeycolour}{\ealtext{#1}}}

% One translated word above one English word. Built from kernel
% primitives, never a tabular, which would break wherever the sheet
% loads array, booktabs or tabularx. The box is as wide as the wider of
% the two, so a long translation pushes its neighbours aside instead of
% overlapping them, and it claims its full height so a table row or a
% TikZ node makes room for it.
\newcommand{\ealgloss}[2]{%
  \sbox{\ealboxen}{\ealkey{#1}}%
  \sbox{\ealboxtr}{\small\textcolor{ealtwocolour}{\ealtext{#2}}}%
  \setlength{\ealglosswd}{\wd\ealboxen}%
  \ifdim\wd\ealboxtr>\ealglosswd\setlength{\ealglosswd}{\wd\ealboxtr}\fi
  \leavevmode
  \makebox[\ealglosswd][c]{%
    \makebox[0pt][c]{\usebox{\ealboxen}}%
    \makebox[0pt][c]{\raisebox{\ealglossraise}{\usebox{\ealboxtr}}}%
  }%
}

% opens up line spacing around a block of glosses, so the raised
% translations sit clear of the line above
\newenvironment{ealglossed}{\par\linespread{2.1}\selectfont\ignorespaces}{\par}

% a whole translated sentence above its English counterpart. block level,
% so it does not belong inside a tabular cell
\newcommand{\ealpara}[2]{%
  \par\addvspace{0.5\baselineskip}%
  {\color{ealtwocolour}\ealtextblock{#1}}%
  \penalty200
  {\color{ealenglish}#2\par}%
  \addvspace{0.5\baselineskip}%
}

\begin{document}

\frame{\titlepage}

\begin{frame}
  \vfill
  \begin{center}
    \ealpara{\ealkeytr{فیصد}: فوری سوالات}{\ealkey{Percentages}: Quick Questions}
    \vspace{0.5em}
    \ealpara{منی وائٹ بورڈ سوالات، حل کے ساتھ}{Mini whiteboard questions, with solutions}
  \end{center}
  \vfill
\end{frame}

\begin{frame}
  \ealpara{فہرست}{Contents}
  \small
  \tableofcontents
\end{frame}

%================================================================
\section{Converting between decimals and percentages}
\begin{frame}[plain]
  \vfill
  \begin{center}
    {\usebeamercolor[fg]{structure}\Huge\bfseries
     \ealpara{\ealkeytr{اعشاریہ} اور \ealkeytr{فیصد} کے درمیان تبدیل کرنا}{Converting between \ealkey{decimals} and \ealkey{percentages}}}
  \end{center}
  \vfill
\end{frame}

\qq[\LARGE]{\ealpara{لکھیں $0.5$ کو \ealkeytr{فیصد} کے طور پر}{Write $0.5$ as a \ealkey{percentage}}}{$0.5\times 100=50\%$}
\qq[\LARGE]{\ealpara{لکھیں $0.25$ کو \ealkeytr{فیصد} کے طور پر}{Write $0.25$ as a \ealkey{percentage}}}{$0.25\times 100=25\%$}
\qq[\LARGE]{\ealpara{لکھیں $30\%$ کو \ealkeytr{اعشاریہ} کے طور پر}{Write $30\%$ as a \ealkey{decimal}}}{$30\div 100=0.3$}
\qq[\LARGE]{\ealpara{لکھیں $0.07$ کو \ealkeytr{فیصد} کے طور پر}{Write $0.07$ as a \ealkey{percentage}}}{$0.07\times 100=7\%$}
\qq[\LARGE]{\ealpara{لکھیں $6\%$ کو \ealkeytr{اعشاریہ} کے طور پر}{Write $6\%$ as a \ealkey{decimal}}}{$6\div 100=0.06$}
\qq[\LARGE]{\ealpara{لکھیں $0.125$ کو \ealkeytr{فیصد} کے طور پر}{Write $0.125$ as a \ealkey{percentage}}}{$0.125\times 100=12.5\%$}
\qq[\LARGE]{\ealpara{لکھیں $4.5\%$ کو \ealkeytr{اعشاریہ} کے طور پر}{Write $4.5\%$ as a \ealkey{decimal}}}{$4.5\div 100=0.045$}
\qq[\LARGE]{\ealpara{لکھیں $1.3$ کو \ealkeytr{فیصد} کے طور پر}{Write $1.3$ as a \ealkey{percentage}}}{$1.3\times 100=130\%$}
\qq[\LARGE]{\ealpara{لکھیں $250\%$ کو \ealkeytr{اعشاریہ} کے طور پر}{Write $250\%$ as a \ealkey{decimal}}}{$250\div 100=2.5$}
\qq[\LARGE]{\ealpara{لکھیں $0.0025$ کو \ealkeytr{فیصد} کے طور پر}{Write $0.0025$ as a \ealkey{percentage}}}{$0.0025\times 100=0.25\%$}

%================================================================
\section{Writing one number as a percentage of another}
\begin{frame}[plain]
  \vfill
  \begin{center}
    {\usebeamercolor[fg]{structure}\Huge\bfseries
     \ealpara{ایک عدد کو دوسرے کے \ealkeytr{فیصد} کے طور پر لکھنا}{Writing one number as a \ealkey{percentage} of another}}
  \end{center}
  \vfill
\end{frame}

\qq[\LARGE][\Large]{\ealpara{لکھیں $25$ کو $100$ کے \ealkeytr{فیصد} کے طور پر}{Write $25$ as a \ealkey{percentage} of $100$}}{$\dfrac{25}{100}=25\%$}
\qq[\LARGE][\Large]{\ealpara{لکھیں $7$ کو $10$ کے \ealkeytr{فیصد} کے طور پر}{Write $7$ as a \ealkey{percentage} of $10$}}{$\dfrac{7}{10}=\dfrac{70}{100}=70\%$}
\qq[\LARGE][\Large]{\ealpara{لکھیں $13$ کو $20$ کے \ealkeytr{فیصد} کے طور پر}{Write $13$ as a \ealkey{percentage} of $20$}}{$\dfrac{13}{20}=\dfrac{65}{100}=65\%$}
\qq[\LARGE][\Large]{\ealpara{لکھیں $12$ کو $25$ کے \ealkeytr{فیصد} کے طور پر}{Write $12$ as a \ealkey{percentage} of $25$}}{$\dfrac{12}{25}=\dfrac{48}{100}=48\%$}
\qq[\LARGE][\Large]{\ealpara{لکھیں $18$ کو $40$ کے \ealkeytr{فیصد} کے طور پر}{Write $18$ as a \ealkey{percentage} of $40$}}{$\dfrac{18}{40}=\dfrac{45}{100}=45\%$}
\qq[\LARGE][\Large]{\ealpara{لکھیں $21$ کو $50$ کے \ealkeytr{فیصد} کے طور پر}{Write $21$ as a \ealkey{percentage} of $50$}}{$\dfrac{21}{50}=\dfrac{42}{100}=42\%$}
\qq[\LARGE][\Large]{\ealpara{لکھیں $9$ کو $12$ کے \ealkeytr{فیصد} کے طور پر}{Write $9$ as a \ealkey{percentage} of $12$}}{$\dfrac{9}{12}=\dfrac{3}{4}=\dfrac{75}{100}=75\%$}
\qq[\LARGE][\Large]{\ealpara{لکھیں $27$ کو $60$ کے \ealkeytr{فیصد} کے طور پر}{Write $27$ as a \ealkey{percentage} of $60$}}{$\dfrac{27}{60}=\dfrac{9}{20}=\dfrac{45}{100}=45\%$}
\qq[\LARGE][\Large]{\ealpara{لکھیں $34$ کو $80$ کے \ealkeytr{فیصد} کے طور پر}{Write $34$ as a \ealkey{percentage} of $80$}}{$\dfrac{34}{80}=\dfrac{17}{40}=\dfrac{42.5}{100}=42.5\%$}
\qq[\LARGE][\Large]{\ealpara{ایک ٹیسٹ میں علی نے $75$ میں سے $57$ اسکور کیے۔ علی نے کتنا \ealkeytr{فیصد} اسکور کیا؟}{In a test Ali scored $57$ out of $75$. What \ealkey{percentage} did Ali score?}}{$\dfrac{57}{75}=\dfrac{19}{25}=\dfrac{76}{100}=76\%$}

%================================================================
\section{Basic percentages without a calculator}
\begin{frame}[plain]
  \vfill
  \begin{center}
    {\usebeamercolor[fg]{structure}\Huge\bfseries
     \ealpara{کیلکولیٹر کے بغیر بنیادی \ealkeytr{فیصد}}{Basic \ealkey{percentages} without a calculator}}
  \end{center}
  \vfill
\end{frame}

\qq{\ealpara{$80$ کا $50\%$}{$50\%$ of $80$}}{$80\div 2=40$}
\qq{\ealpara{$70$ کا $10\%$}{$10\%$ of $70$}}{$70\div 10=7$}
\qq{\ealpara{$45$ کا $10\%$}{$10\%$ of $45$}}{$45\div 10=4.5$}
\qq{\ealpara{$300$ کا $1\%$}{$1\%$ of $300$}}{$300\div 100=3$}
\qq{\ealpara{$45$ کا $20\%$}{$20\%$ of $45$}}{\steps{$10\%$ & $4.5$ \\ $20\%$ & $9$}}
\qq{\ealpara{$80$ کا $5\%$}{$5\%$ of $80$}}{\steps{$10\%$ & $8$ \\ $5\%$ & $4$}}
\qq{\ealpara{$250$ کا $30\%$}{$30\%$ of $250$}}{\steps{$10\%$ & $25$ \\ $30\%$ & $75$}}
\qq{\ealpara{$42$ کا $1\%$}{$1\%$ of $42$}}{$42\div 100=0.42$}
\qq{\ealpara{$60$ کا $15\%$}{$15\%$ of $60$}}{\steps{$10\%$ & $6$ \\ $5\%$ & $3$ \\ $15\%$ & $9$}}
\qq{\ealpara{$240$ کا $35\%$}{$35\%$ of $240$}}{\steps{$10\%$ & $24$ \\ $30\%$ & $72$ \\ $5\%$ & $12$ \\ $35\%$ & $84$}}

%================================================================
\section{Increasing and decreasing without a calculator}
\begin{frame}[plain]
  \vfill
  \begin{center}
    {\usebeamercolor[fg]{structure}\Huge\bfseries
     \ealpara{کیلکولیٹر کے بغیر \ealkeytr{اضافہ کرنا} اور \ealkeytr{کم کرنا}}{\ealkey{Increasing} and \ealkey{decreasing} without a calculator}}
  \end{center}
  \vfill
\end{frame}

\qq[\LARGE]{\ealpara{$80$ میں $10\%$ کا \ealkeytr{اضافہ کرنا}}{\ealkey{Increase} $80$ by $10\%$}}{\steps{$10\%$ & $8$}\par\vspace{0.4em}$80+8=88$}
\qq[\LARGE]{\ealpara{$60$ میں $50\%$ کی \ealkeytr{کم کرنا}}{\ealkey{Decrease} $60$ by $50\%$}}{\steps{$50\%$ & $30$}\par\vspace{0.4em}$60-30=30$}
\qq[\LARGE]{\ealpara{$40$ میں $25\%$ کا \ealkeytr{اضافہ کرنا}}{\ealkey{Increase} $40$ by $25\%$}}{\steps{$25\%$ & $10$}\par\vspace{0.4em}$40+10=50$}
\qq[\LARGE]{\ealpara{$90$ میں $10\%$ کی \ealkeytr{کم کرنا}}{\ealkey{Decrease} $90$ by $10\%$}}{\steps{$10\%$ & $9$}\par\vspace{0.4em}$90-9=81$}
\qq[\LARGE]{\ealpara{\pounds $250$ میں $20\%$ کا \ealkeytr{اضافہ کرنا}}{\ealkey{Increase} \pounds $250$ by $20\%$}}{\steps{$10\%$ & \pounds $25$ \\ $20\%$ & \pounds $50$}\par\vspace{0.4em} \pounds $250+$ \pounds $50=$ \pounds $300$}
\qq[\LARGE]{\ealpara{$45$ میں $20\%$ کی \ealkeytr{کم کرنا}}{\ealkey{Decrease} $45$ by $20\%$}}{\steps{$10\%$ & $4.5$ \\ $20\%$ & $9$}\par\vspace{0.4em}$45-9=36$}
\qq[\LARGE]{\ealpara{$120$ میں $5\%$ کا \ealkeytr{اضافہ کرنا}}{\ealkey{Increase} $120$ by $5\%$}}{\steps{$10\%$ & $12$ \\ $5\%$ & $6$}\par\vspace{0.4em}$120+6=126$}
\qq[\LARGE]{\ealpara{\pounds $64$ میں $25\%$ کی \ealkeytr{کم کرنا}}{\ealkey{Decrease} \pounds $64$ by $25\%$}}{\steps{$25\%$ & \pounds $16$}\par\vspace{0.4em} \pounds $64-$ \pounds $16=$ \pounds $48$}
\qq[\LARGE]{\ealpara{$350$ میں $30\%$ کا \ealkeytr{اضافہ کرنا}}{\ealkey{Increase} $350$ by $30\%$}}{\steps{$10\%$ & $35$ \\ $30\%$ & $105$}\par\vspace{0.4em}$350+105=455$}
\qq[\LARGE]{\ealpara{$480$ میں $35\%$ کی \ealkeytr{کم کرنا}}{\ealkey{Decrease} $480$ by $35\%$}}{\steps{$10\%$ & $48$ \\ $30\%$ & $144$ \\ $5\%$ & $24$ \\ $35\%$ & $168$}%
   \par\vspace{0.4em}$480-168=312$}

%================================================================
\section{Finding the original value without a calculator}
\begin{frame}[plain]
  \vfill
  \begin{center}
    {\usebeamercolor[fg]{structure}\Huge\bfseries
     \ealpara{کیلکولیٹر کے بغیر \ealkeytr{اصل مقدار} معلوم کرنا}{Finding the \ealkey{original value} without a calculator}}
  \end{center}
  \vfill
\end{frame}

\qq[\Large][\large]{\ealpara{ایک عدد کا $10\%$، $7$ ہے۔ وہ عدد کیا ہے؟}{$10\%$ of a number is $7$. What is the number?}}{\ealpara{$10\%$، $7$ ہے؛ $100\%$، ؟ ہے}{$10\%$ is $7$; $100\%$ is $?$}\barmodel{10}{1}{$10\%$ is $7$}{$100\%$ is $?$}\par\vspace{0.9em}
   \steps{$10\%$ & $7$ \\ $100\%$ & $70$}}

\qq[\Large][\large]{\ealpara{ایک عدد کا $50\%$، $24$ ہے۔ وہ عدد کیا ہے؟}{$50\%$ of a number is $24$. What is the number?}}{\ealpara{$50\%$، $24$ ہے؛ $100\%$، ؟ ہے}{$50\%$ is $24$; $100\%$ is $?$}\barmodel{2}{1}{$50\%$ is $24$}{$100\%$ is $?$}\par\vspace{0.9em}
   \steps{$50\%$ & $24$ \\ $100\%$ & $48$}}

\qq[\Large][\large]{\ealpara{ایک عدد کا $20\%$، $12$ ہے۔ وہ عدد کیا ہے؟}{$20\%$ of a number is $12$. What is the number?}}{\ealpara{$20\%$، $12$ ہے؛ $100\%$، ؟ ہے}{$20\%$ is $12$; $100\%$ is $?$}\barmodel{10}{2}{$20\%$ is $12$}{$100\%$ is $?$}\par\vspace{0.9em}
   \steps{$20\%$ & $12$ \\ $10\%$ & $6$ \\ $100\%$ & $60$}}

\qq[\Large][\large]{\ealpara{ایک عدد کا $80\%$، $64$ ہے۔ وہ عدد کیا ہے؟}{$80\%$ of a number is $64$. What is the number?}}{\ealpara{$80\%$، $64$ ہے؛ $100\%$، ؟ ہے}{$80\%$ is $64$; $100\%$ is $?$}\barmodel{10}{8}{$80\%$ is $64$}{$100\%$ is $?$}\par\vspace{0.9em}
   \steps{$80\%$ & $64$ \\ $10\%$ & $8$ \\ $100\%$ & $80$}}

\qq[\Large][\large]{\ealpara{ایک عدد کا $30\%$، $18$ ہے۔ وہ عدد کیا ہے؟}{$30\%$ of a number is $18$. What is the number?}}{\ealpara{$30\%$، $18$ ہے؛ $100\%$، ؟ ہے}{$30\%$ is $18$; $100\%$ is $?$}\barmodel{10}{3}{$30\%$ is $18$}{$100\%$ is $?$}\par\vspace{0.9em}
   \steps{$30\%$ & $18$ \\ $10\%$ & $6$ \\ $100\%$ & $60$}}

\qq[\Large][\large]{\ealpara{ایک عدد کا $25\%$، $15$ ہے۔ وہ عدد کیا ہے؟}{$25\%$ of a number is $15$. What is the number?}}{\ealpara{$25\%$، $15$ ہے؛ $100\%$، ؟ ہے}{$25\%$ is $15$; $100\%$ is $?$}\barmodel{4}{1}{$25\%$ is $15$}{$100\%$ is $?$}\par\vspace{0.9em}
   \steps{$25\%$ & $15$ \\ $100\%$ & $60$}}

\qq[\Large][\large]{\ealpara{ایک عدد کا $15\%$، $12$ ہے۔ وہ عدد کیا ہے؟}{$15\%$ of a number is $12$. What is the number?}}{\ealpara{$15\%$، $12$ ہے؛ $100\%$، ؟ ہے}{$15\%$ is $12$; $100\%$ is $?$}\barmodel{20}{3}{$15\%$ is $12$}{$100\%$ is $?$}\par\vspace{0.9em}
   \steps{$15\%$ & $12$ \\ $5\%$ & $4$ \\ $100\%$ & $80$}}

\qq[\Large][\large]{\ealpara{ایک کوٹ سیل میں $20\%$ سے \ealkeytr{کم کیا گیا ہے}۔ اب اس کی قیمت \pounds $48$ ہے۔ اصل قیمت کیا تھی؟}{A coat is \ealkey{reduced} by $20\%$ in a sale. It now costs \pounds $48$. What was the original price?}}{\ealpara{$80\%$، \pounds $48$ ہے؛ $100\%$، ؟ ہے}{$80\%$ is \pounds $48$; $100\%$ is $?$}\barmodel{10}{8}{$80\%$ is \pounds $48$}{$100\%$ is $?$}\par\vspace{0.9em}
   \steps{$80\%$ & \pounds $48$ \\ $10\%$ & \pounds $6$ \\ $100\%$ & \pounds $60$}}

\qq[\Large][\large]{\ealpara{ایک قیمت $10\%$ بڑھ کر \pounds $88$ ہو جاتی ہے۔ بڑھنے سے پہلے قیمت کیا تھی؟}{A price rises by $10\%$ to \pounds $88$. What was the price before the rise?}}{\ealpara{$100\%$، ؟ ہے؛ $110\%$، \pounds $88$ ہے}{$100\%$ is $?$; $110\%$ is \pounds $88$}\barmodel{11}{10}{$100\%$ is $?$}{$110\%$ is \pounds $88$}\par\vspace{0.9em}
   \steps{$110\%$ & \pounds $88$ \\ $10\%$ & \pounds $8$ \\ $100\%$ & \pounds $80$}}

\qq[\Large][\large]{\ealpara{ایک گھر کی قیمت $30\%$ بڑھ کر \pounds $169\,000$ ہو جاتی ہے۔ بڑھنے سے پہلے اس کی قیمت کتنی تھی؟}{A house rises in value by $30\%$ to \pounds $169\,000$. What was it worth before the rise?}}{\ealpara{$100\%$، ؟ ہے؛ $130\%$، \pounds $169\,000$ ہے}{$100\%$ is $?$; $130\%$ is \pounds $169\,000$}\barmodel{13}{10}{$100\%$ is $?$}{$130\%$ is \pounds $169\,000$}\par\vspace{0.9em}
   \steps{$130\%$ & \pounds $169\,000$ \\ $10\%$ & \pounds $13\,000$ \\ $100\%$ & \pounds $130\,000$}}

%================================================================
\section{Percentage word problems}
\begin{frame}[plain]
  \vfill
  \begin{center}
    {\usebeamercolor[fg]{structure}\Huge\bfseries
     \ealpara{\ealkeytr{فیصد} کے لفظی مسائل}{\ealkey{Percentage} word problems}}
  \end{center}
  \vfill
\end{frame}

\qq[\Large][\Large]{\ealpara{ایک کوٹ کی قیمت \pounds $40$ ہے۔ سیل میں اسے $10\%$ سے \ealkeytr{کم کیا گیا}۔ سیل قیمت کیا ہے؟}{A coat costs \pounds $40$. In a sale it is \ealkey{reduced} by $10\%$. What is the sale price?}}{\steps{$10\%$ & \pounds $4$}\par\vspace{0.4em}
   \pounds $40-$ \pounds $4=$ \pounds $36$}

\qq[\Large][\Large]{\ealpara{ایک کلاس کے $30$ طلباء میں سے $30\%$ پیدل اسکول جاتے ہیں۔ کتنے طلباء پیدل اسکول جاتے ہیں؟}{$30\%$ of a class of $30$ students walk to school. How many students walk to school?}}{\ealpara{$9$ طلباء}{$9$ students}\steps{$10\%$ & $3$ \\ $30\%$ & $9$}\par\vspace{0.4em}
   $9$ students}

\qq[\Large][\Large]{\ealpara{ایک فون کی قیمت \pounds $250$ کے علاوہ $20\%$ VAT ہے۔ کل قیمت کیا ہے؟}{A phone costs \pounds $250$ plus $20\%$ VAT. What is the total price?}}{\steps{$10\%$ & \pounds $25$ \\ $20\%$ & \pounds $50$}\par\vspace{0.4em}
   \pounds $250+$ \pounds $50=$ \pounds $300$}

\qq[\Large][\Large]{\ealpara{سیم ماہانہ \pounds $800$ کماتا ہے اور اسے $5\%$ تنخواہ میں اضافہ دیا جاتا ہے۔ سیم کی نئی ماہانہ تنخواہ کیا ہے؟}{Sam earns \pounds $800$ a month and is given a $5\%$ pay rise. What is Sam's new monthly pay?}}{\steps{$10\%$ & \pounds $80$ \\ $5\%$ & \pounds $40$}\par\vspace{0.4em}
   \pounds $800+$ \pounds $40=$ \pounds $840$}

\qq[\Large][\Large]{\ealpara{ایک جیکٹ کی قیمت \pounds $60$ تھی اور اب \pounds $45$ ہے۔ \ealkeytr{فیصد} \ealkeytr{کمی} کیا ہے؟}{A jacket cost \pounds $60$ and now costs \pounds $45$. What is the \ealkey{percentage} \ealkey{reduction}?}}{\ealpara{\ealkeytr{کمی} $=$ \pounds $15$ \par\vspace{0.4em}
   $\dfrac{15}{60}=\dfrac{25}{100}=25\%$}{\ealkey{Reduction} $=$ \pounds $15$ \par\vspace{0.4em}
   $\dfrac{15}{60}=\dfrac{25}{100}=25\%$}}

\qq[\Large][\Large]{\ealpara{ایک قصبے کی آبادی $8000$ سے $8600$ ہو جاتی ہے۔ \ealkeytr{فیصد} \ealkeytr{اضافہ} کیا ہے؟}{A town's population grows from $8000$ to $8600$. What is the \ealkey{percentage} \ealkey{increase}?}}{\ealpara{\ealkeytr{اضافہ} $=600$ \par\vspace{0.4em}
   $\dfrac{600}{8000}=\dfrac{7.5}{100}=7.5\%$}{\ealkey{Increase} $=600$ \par\vspace{0.4em}
   $\dfrac{600}{8000}=\dfrac{7.5}{100}=7.5\%$}}

\qq[\Large][\large]{\ealpara{ایک لیپ ٹاپ سیل میں \pounds $480$ کا ہے، $20\%$ \ealkeytr{رعایت} کے بعد۔ سیل سے پہلے قیمت کیا تھی؟}{A laptop costs \pounds $480$ in a sale, after a $20\%$ \ealkey{discount}. What was the price before the sale?}}{\ealpara{$80\%$، \pounds $480$ ہے؛ $100\%$، ؟ ہے}{$80\%$ is \pounds $480$; $100\%$ is $?$}\barmodel{10}{8}{$80\%$ is \pounds $480$}{$100\%$ is $?$}\par\vspace{0.9em}
   \steps{$80\%$ & \pounds $480$ \\ $10\%$ & \pounds $60$ \\ $100\%$ & \pounds $600$}}

\qq[\Large][\Large]{\ealpara{مایا نے ایک ٹیسٹ میں $40$ میں سے $34$ اسکور کیے۔ بین نے $82\%$ اسکور کیے۔ کس نے بہتر کارکردگی دکھائی؟}{Maya scored $34$ out of $40$ in a test. Ben scored $82\%$. Who did better?}}{\ealpara{$\dfrac{34}{40}=\dfrac{85}{100}=85\%$ \par\vspace{0.4em}
   $85\%>82\%$، لہٰذا مایا نے بہتر کارکردگی دکھائی}{$\dfrac{34}{40}=\dfrac{85}{100}=85\%$ \par\vspace{0.4em}
   $85\%>82\%$, so Maya did better}}

\qq[\Large][\Large]{\ealpara{ایک ٹیلی ویژن کی قیمت \pounds $1200$ ہے۔ اسے $15\%$ سے \ealkeytr{کم کیا گیا}، پھر نئی قیمت کے $10\%$ سے مزید کم کیا گیا۔ حتمی قیمت کیا ہے؟}{A television costs \pounds $1200$. It is \ealkey{reduced} by $15\%$, then by a further $10\%$ of the new price. What is the final price?}}{\ealpara{$15\%$ از \pounds $1200$، \pounds $180$؛ نئی قیمت، \pounds $1020$؛ $10\%$ از \pounds $1020$، \pounds $102$۔ حتمی قیمت $=$ \pounds $918$}{$15\%$ of \pounds $1200$, \pounds $180$; new price, \pounds $1020$; $10\%$ of \pounds $1020$, \pounds $102$; Final price $=$ \pounds $918$}
\steps{$15\%$ of \pounds $1200$ & \pounds $180$ \\
          new price & \pounds $1020$ \\
          $10\%$ of \pounds $1020$ & \pounds $102$}\par\vspace{0.4em}
   Final price $=$ \pounds $918$}

\qq[\Large][\Large]{\ealpara{ایک کار کی قیمت \pounds $12\,000$ ہے۔ یہ پہلے سال میں اپنی قیمت کا $20\%$ کھو دیتی ہے اور دوسرے سال میں $15\%$ کھو دیتی ہے۔ دو سال بعد اس کی قیمت کتنی ہوگی؟}{A car costs \pounds $12\,000$. It loses $20\%$ of its value in the first year and $15\%$ of its value in the second year. What is it worth after two years?}}{\ealpara{پہلے سال کے بعد \pounds $12\,000\times 0.8=$ \pounds $9600$؛ دوسرے سال کے بعد \pounds $9600\times 0.85=$ \pounds $8160$}{after year $1$: \pounds $12\,000\times 0.8=$ \pounds $9600$; after year $2$: \pounds $9600\times 0.85=$ \pounds $8160$}
\steps{after year $1$ & \pounds $12\,000\times 0.8=$ \pounds $9600$ \\
          after year $2$ & \pounds $9600\times 0.85=$ \pounds $8160$}}

%================================================================
\section{Harder percentages of a number}
\begin{frame}[plain]
  \vfill
  \begin{center}
    {\usebeamercolor[fg]{structure}\Huge\bfseries
     \ealpara{عدد کے مشکل \ealkeytr{فیصد}}{Harder \ealkey{percentages} of a number}}
  \end{center}
  \vfill
\end{frame}

\qq{\ealpara{$200$ کا $11\%$}{$11\%$ of $200$}}{\steps{$10\%$ & $20$ \\ $1\%$ & $2$ \\ $11\%$ & $22$}}
\qq{\ealpara{$400$ کا $21\%$}{$21\%$ of $400$}}{\steps{$10\%$ & $40$ \\ $20\%$ & $80$ \\ $1\%$ & $4$ \\ $21\%$ & $84$}}
\qq{\ealpara{$40$ کا $17.5\%$}{$17.5\%$ of $40$}}{\steps{$10\%$ & $4$ \\ $5\%$ & $2$ \\ $2.5\%$ & $1$ \\ $17.5\%$ & $7$}}
\qq{\ealpara{$80$ کا $12.5\%$}{$12.5\%$ of $80$}}{\steps{$10\%$ & $8$ \\ $2.5\%$ & $2$ \\ $12.5\%$ & $10$}}
\qq{\ealpara{$240$ کا $17.5\%$}{$17.5\%$ of $240$}}{\steps{$10\%$ & $24$ \\ $5\%$ & $12$ \\ $2.5\%$ & $6$ \\ $17.5\%$ & $42$}}
\qq{\ealpara{$150$ کا $32\%$}{$32\%$ of $150$}}{\steps{$10\%$ & $15$ \\ $30\%$ & $45$ \\ $1\%$ & $1.5$ \\ $2\%$ & $3$ \\ $32\%$ & $48$}}
\qq{\ealpara{$68$ کا $17.5\%$}{$17.5\%$ of $68$}}{\steps{$10\%$ & $6.8$ \\ $5\%$ & $3.4$ \\ $2.5\%$ & $1.7$ \\ $17.5\%$ & $11.9$}}
\qq{\ealpara{$96$ کا $87.5\%$}{$87.5\%$ of $96$}}{\steps{$100\%$ & $96$ \\ $10\%$ & $9.6$ \\ $2.5\%$ & $2.4$ \\ $12.5\%$ & $12$ \\ $87.5\%$ & $96-12=84$}}
\qq{\ealpara{$400$ کا $6.5\%$}{$6.5\%$ of $400$}}{\steps{$1\%$ & $4$ \\ $6\%$ & $24$ \\ $0.5\%$ & $2$ \\ $6.5\%$ & $26$}}
\qq[\Huge][\LARGE]{\ealpara{\pounds $312$ کا $17.5\%$}{$17.5\%$ of \pounds $312$}}{\steps{$10\%$ & \pounds $31.20$ \\ $5\%$ & \pounds $15.60$ \\
          $2.5\%$ & \pounds $7.80$ \\ $17.5\%$ & \pounds $54.60$}}

\end{document}